\section{Zadanie 2}
% opis
Porównujemy algorytmy rozwiązujące macierze dolnotrójkątne
\textbf{Działanie SolveL}
\begin{enumerate}
    \item pętla $i$ po indeksach $2$ do długości $x$
    \item pętla po indeksach 1 do $i$
    \item zliczanie - mnożymy po kolei elementy 
\end{enumerate}
\textbf{Działanie SolveL1}
\begin{enumerate}
    \item pętla po indeksach wierszy, 
    \item zliczanie - mnożenie macierzowe 
\end{enumerate}

% Please add the following required packages to your document preamble:
% \usepackage{graphicx}

\begin{center}
% Please add the following required packages to your document preamble:
% \usepackage{graphicx}
\resizebox{0.7\textwidth}{!}{%
\begin{tabular}{|l|l|l|l|l|}
\hline
$n$          & SolveL     & błąd          & SolveL1       & błąd          \\ \hline
1000       & $1.43726$s & $6.2607e+159$ & $0.00959992$s & $1.2082e+159$ \\ \hline
10000      & $135.994$s & $1.9449e+295$ & $0.278352$s   & $1.8671e+295$ \\ \hline
1000 + eye & $1.36324$s & $4.7684e-06$  & $0.00939488$s & $6.4969e-06$  \\ \hline
\end{tabular}%
}
\end{center}


% wnioski 
\noindent
Widzimy że rozwiązywanie układów dolnotrójkątnych jest bardzo szybkie. Trzeba zwrócić uwagę na błąd pomiarowy. Bierze się to z iloczyn złożony z wielu małych liczb. Takie wyrażenie daży do $0$. Kiedy dzielimy przez odwrotność, utrzymujemy bardzo duże błędy.